% chapters/01/section-02.tex
\section{Buku Ajar dan RPS Berbasis OBE}
\label{sec:obe}

Buku ini disusun mengikuti Rencana Pembelajaran Semester (RPS) mata kuliah
\NamaMK\ yang menganut pendekatan \istilah{Outcome-Based Education} (OBE).
Bagian ini menjelaskan apa arti pendekatan tersebut, bagaimana buku ini
menerjemahkannya menjadi urutan bab, dan apa yang perlu Anda ketahui agar
dapat membaca buku ini sebagaimana dimaksudkan penulisnya.

\subsection{Prinsip \istilah{Outcome-Based Education}}

\istilah{Outcome-Based Education} adalah pendekatan penyusunan pembelajaran
yang menetapkan lebih dahulu apa yang harus mampu dilakukan mahasiswa,
kemudian memilih materi, aktivitas, dan asesmen yang mengarah pada kemampuan
tersebut. Urutan berpikirnya terbalik dibandingkan penyusunan mata kuliah
tradisional yang sering berangkat dari daftar materi yang ingin diajarkan.

Dua konsekuensi pendekatan ini penting bagi Anda sebagai mahasiswa.

\begin{enumerate}
  \item \textbf{Penilaian berbasis kriteria, bukan peringkat.} Nilai Anda
        ditentukan oleh sejauh mana kriteria ketercapaian terpenuhi, bukan
        oleh posisi Anda dibandingkan teman sekelas. Jika seluruh mahasiswa
        memenuhi kriteria, seluruhnya memperoleh nilai yang baik; tidak ada
        kuota nilai rendah yang harus dipenuhi.
  \item \textbf{Capaian diumumkan sejak awal.} Tidak ada maksud tersembunyi di
        balik penilaian. Setiap bab pada buku ini memuat Sub-CPMK, indikator
        pencapaian, aktivitas yang harus dikerjakan, dan instrumen penilaian
        yang dipakai. Anda berhak mengetahui ukuran keberhasilan sebelum
        pekerjaan dinilai.
\end{enumerate}

Pendekatan ini juga diikuti oleh banyak program studi ilmu komputer di
Indonesia yang merancang mata kuliah pemrograman permainan dengan cara
serupa, walaupun cakupan dan urutan materinya berbeda-beda
\cite{its-rps-game,unima-rps-game,polibatam-kurikulum}. Perbedaan tersebut
wajar karena bergantung pada capaian pembelajaran lulusan yang dibebankan
pada masing-masing program studi.

\subsection{Penyelarasan Konstruktif}

Pelaksanaan OBE menuntut apa yang disebut \istilah{constructive alignment},
yaitu keadaan ketika capaian, materi, aktivitas belajar, dan asesmen menunjuk
arah yang sama. Gambar~\ref{fig:alur-obe} memperlihatkan alur penyusunan
yang dipakai pada buku ini.

\begin{figure}[htbp]
\centering
% chapters/01/figures/alur-obe.tikz
% Penyelarasan konstruktif: dari capaian lulusan sampai bukti belajar.
\begin{tikzpicture}[
  kotak/.style={
    rectangle, rounded corners=2pt, draw=ubPrimary, thick, fill=ubAlt,
    align=center, inner sep=5pt, minimum height=10mm, text width=32mm,
    font=\footnotesize
  },
  sorot/.style={kotak, draw=ubAccent, fill=white},
  bukti/.style={kotak, draw=ubGreen, fill=ubGreenBg},
  panah/.style={-{Stealth[length=2.4mm]}, thick, draw=ubAccent},
  kecil/.style={font=\scriptsize\itshape\color{ubPrimary}, align=center}
]

\node[kotak] (cpl) at (0,0) {\textbf{CPL03} dan \textbf{CPL04}\\[1pt]\scriptsize capaian pembelajaran lulusan};
\node[kotak, below=10mm of cpl] (cpmk) {\textbf{CPMK1}--\textbf{CPMK4}\\[1pt]\scriptsize capaian pembelajaran mata kuliah};
\node[kotak, below=10mm of cpmk] (sub) {\textbf{13 Sub-CPMK}\\[1pt]\scriptsize capaian pembelajaran tiap bab};

\node[sorot] (materi) at (-46mm,-72mm) {\textbf{Materi}\\[1pt]\scriptsize bahan kajian pada bab};
\node[sorot] (aktivitas) at (0,-72mm) {\textbf{Aktivitas}\\[1pt]\scriptsize praktikum dan proyek};
\node[sorot] (asesmen) at (46mm,-72mm) {\textbf{Asesmen}\\[1pt]\scriptsize rubrik dan instrumen};

\node[bukti] (bukti) at (0,-100mm) {\textbf{Bukti belajar}\\[1pt]\scriptsize kode, level, laporan, dan refleksi};

\draw[panah] (cpl) -- (cpmk);
\draw[panah] (cpmk) -- (sub);
\draw[panah] (sub.south) -- ++(0,-6mm) -| (materi.north);
\draw[panah] (sub.south) -- (aktivitas.north);
\draw[panah] (sub.south) -- ++(0,-6mm) -| (asesmen.north);
\draw[panah] (materi.south) -- ++(0,-6mm) -| (bukti.north);
\draw[panah] (aktivitas) -- (bukti);
\draw[panah] (asesmen.south) -- ++(0,-6mm) -| (bukti.north);

\node[kecil, right=4mm of sub, text width=42mm, anchor=west]
  {capaian ditetapkan lebih dahulu; materi, aktivitas, dan asesmen mengikutinya};
\end{tikzpicture}

\caption{Alur penyelarasan konstruktif dari capaian lulusan sampai bukti belajar}
\label{fig:alur-obe}
\end{figure}

Alur tersebut dapat dibaca sebagai berikut.

\begin{langkah}
  \item \textbf{CPL} ditetapkan pada tingkat program studi. Untuk mata kuliah
        ini, dua capaian yang dibebankan adalah \CPL{03} tentang perancangan
        dan analisis algoritma serta penerapannya pada bahasa pemrograman
        tertentu, dan \CPL{04} tentang perancangan, implementasi, serta
        evaluasi solusi berbasis \istilah{computing}.
  \item \textbf{\CPMK} menjabarkan CPL menjadi empat kemampuan yang dapat
        diuji pada akhir mata kuliah, sebagaimana ditampilkan pada
        Tabel~\ref{tab:capaian-bukti}.
  \item \textbf{Sub-CPMK} memecah setiap CPMK menjadi kemampuan yang lebih
        spesifik dan terukur pada tingkat bab. Mata kuliah ini memiliki
        tiga belas Sub-CPMK, dan setiap bab inti pada buku ini menopang satu
        di antaranya.
  \item \textbf{Materi} yang dipelajari dipilih agar kemampuan tersebut dapat
        terbentuk, bukan karena materi itu menarik atau mudah diajarkan.
  \item \textbf{Aktivitas} dirancang agar mahasiswa benar-benar berlatih
        menggunakan kemampuan itu, bukan hanya mendengarnya dijelaskan.
  \item \textbf{Asesmen} disusun mengukur kemampuan yang sama, dengan
        indikator yang tampak pada pekerjaan mahasiswa.
  \item \textbf{Bukti belajar} menjadi akhir alurnya: kode program, level
        yang dapat dimainkan, laporan \istilah{playtest}, dan lembar refleksi.
\end{langkah}

Sebagai contoh, \SUBCPMK{1.2} menuntut kemampuan merancang dan
mengimplementasikan gerak objek berbasis \istilah{delta time} dan vektor pada
Pygame. Materi penopangnya adalah Bab~V. Aktivitasnya adalah praktikum
menggerakkan karakter dengan kecepatan tetap pada 30 FPS dan 60 FPS.
Asesmennya adalah penilaian artefak kode tersebut dengan rubrik, ditambah
butir kuis yang menanyakan mengapa kecepatan harus dikalikan
\istilah{delta time}. Ketiga unsur itu menunjuk kemampuan yang sama, yaitu
\SUBCPMK{1.2} itu sendiri. Contoh pembacaan yang lebih rinci diberikan pada
Bagian~\ref{sec:contoh-bab}.

Penyelarasan semacam ini mencegah dua kesalahan yang sering terjadi pada
mata kuliah pemrograman: mengajarkan banyak hal yang tidak pernah dinilai,
dan menilai sesuatu yang tidak pernah dilatihkan.

\subsection{Keterkaitan Buku Ajar dengan RPS}

Buku ini bukan pengganti RPS, melainkan bahan belajar yang mengikuti RPS
tersebut. Tiga hal perlu Anda ketahui tentang keterkaitannya.

\begin{itemize}
  \item \textbf{Satu bab setara dengan satu materi pembelajaran.} Bab III
        sampai Bab XVI pada buku ini berpasangan satu-satu dengan Materi 1
        sampai Materi 14 pada RPS. Bab I, Bab II, dan Bab XVII bersifat
        pengantar dan penutup sehingga tidak memiliki pasangan materi.
  \item \textbf{Nomor rujukan pustaka sama.} Rujukan pada buku ini bernomor
        \mbox{[R1]} sampai \mbox{[R35]}, persis seperti penomoran pada RPS.
        Jika dosen Anda menyebut rujukan \mbox{[R8]} di kelas, rujukan itu
        menunjuk pustaka yang sama pada kedua dokumen.
  \item \textbf{Peta pemetaan tersedia di bagian awal buku.} Tabel yang
        menghubungkan bab, materi, Sub-CPMK, dan minggu perkuliahan diletakkan
        pada bagian \istilah{Capaian Pembelajaran dan Peta Kompetensi}, bukan
        diulang pada setiap bab. Gunakan tabel tersebut setiap kali Anda
        ingin mengetahui mengapa sebuah bab dibahas pada minggu tertentu.
\end{itemize}

Apabila dosen pengampu menyesuaikan urutan atau penekanan materi sesuai
kebutuhan kelas, urutan bab pada buku ini tetap dapat diikuti karena setiap
bab dirancang sebagai kesatuan yang utuh: memiliki capaian, materi,
praktikum, dan asesmennya sendiri.

\subsection{Peran Setiap Dokumen Rujukan}

Selain RPS dan buku ini, tersedia beberapa kelompok dokumen lain. Membedakan
perannya akan menghemat waktu Anda, terutama ketika mencari jawaban atas
persoalan teknis.

\begin{table}[htbp]
\centering
\caption{Peran setiap dokumen rujukan}
\label{tab:peran-dokumen}
{\setlength{\tabcolsep}{4pt}\footnotesize%
\begin{tabularx}{\textwidth}{@{}L{3.6cm} L{4.6cm} Y@{}}
\toprule
\headrow \textbf{Dokumen} & \textbf{Fungsi} & \textbf{Dipakai ketika} \\
\midrule
RPS mata kuliah &
Rencana resmi: capaian, materi, metode, alokasi waktu, bobot penilaian, dan
rubrik. &
Anda ingin mengetahui kapan suatu materi dibahas, berapa bobotnya, atau apa
kriteria kelulusan. \\
\addlinespace
Buku ajar ini &
Bahan belajar dan lembar kerja: penjelasan materi, contoh, praktikum, latihan,
dan instrumen asesmen per bab. &
Anda sedang mempelajari materi atau mengerjakan praktikum. \\
\addlinespace
Dokumentasi resmi Python, Pygame, dan Unity \cite{pygame-docs,unity-learn} &
Rujukan teknis terkini: daftar fungsi, parameter, dan perilaku pustaka. &
Program Anda berjalan tidak seperti yang diharapkan dan Anda perlu memeriksa
perilaku fungsi yang dipakai. \\
\addlinespace
Buku dan bahan kuliah rujukan \mbox{[R1]}--\mbox{[R11]} &
Penjelasan tambahan, contoh permainan utuh, dan sudut pandang lain atas materi
yang sama. &
Anda ingin memperdalam suatu topik atau melihat contoh penerapan yang lebih
besar. \\
\addlinespace
Repositori Git Anda &
Bukti portofolio: riwayat pekerjaan praktikum mingguan beserta catatannya. &
Sepanjang semester, mulai dari praktikum pertama. \\
\bottomrule
\end{tabularx}}
\end{table}

\begin{catatabox}
Buku ini menjelaskan cara kerja dan alasan di balik setiap teknik, sedangkan
dokumentasi resmi menjelaskan pilihan yang tersedia secara lengkap. Ketika
keduanya berbeda karena pustaka diperbarui, dokumentasi resmi yang berlaku.
Versi perkakas yang dipakai pada buku ini dicatat pada Bab~III untuk Python
dan Pygame, serta pada Bab~XII untuk Unity.
\end{catatabox}
